\chapter{Conclusão}

O desenvolvimento do sistema \textit{SeniorCareManager} demonstrou a relevância da tecnologia como ferramenta de fortalecimento institucional em organizações sociais. A partir da identificação das fragilidades operacionais do Lar dos Velhinhos São Vicente de Paulo e da APAE de Jales, evidenciou-se a necessidade de uma solução que integrasse informações assistenciais, administrativas e logísticas de forma segura e rastreável. O projeto resultou em uma plataforma robusta, composta por módulos independentes e intercomunicáveis, capazes de automatizar processos essenciais, como prontuário eletrônico, controle de estoque, compras, doações e gestão de vacinas.

A utilização de práticas consolidadas de engenharia de software, aliada à modelagem UML, possibilitou a definição clara de requisitos, fluxos e entidades, garantindo coerência estrutural e aderência às necessidades institucionais. Aspectos como acessibilidade, segurança da informação e conformidade com a LGPD reforçaram o compromisso do projeto com a proteção de dados e a inclusão tecnológica.

Assim, o \textit{SeniorCareManager} representa um avanço significativo na modernização das rotinas internas das instituições atendidas, contribuindo para maior eficiência operacional, qualidade assistencial e otimização de recursos. O projeto reafirma o potencial transformador da extensão acadêmica, ao integrar conhecimento técnico e impacto social.
